\section{Rastreabilidade}

A cadeia mínima de evidência é:

\begin{center}
Norma $\downarrow$ \texttt{IMP-*} $\downarrow$ \texttt{TEST-*}
$\downarrow$ arquivo de teste $\downarrow$ resultado
\end{center}

Cada caso permanente recebe um identificador estável por camada:

\begin{itemize}
  \item \texttt{TEST-UNIT-*}: lógica isolada;
  \item \texttt{TEST-INT-*}: integração com serviços locais;
  \item \texttt{TEST-RULES-*}: autorização em Firestore ou Storage Rules;
  \item \texttt{TEST-E2E-*}: fluxo observável no navegador.
\end{itemize}

O sufixo deve ser legível e estável. Quando houver milestone relevante, ele pode
ser incorporado, por exemplo \texttt{TEST-RULES-M9-001}. O registro associa o ID
à referência normativa exata, à feature da matriz, ao arquivo/símbolo exercitado,
ao resultado e à evidência da execução.

\subsection{Exemplo não executado}

O encadeamento abaixo é somente um exemplo de identificação futura; não declara
teste criado nem resultado:

\begin{center}
Seção 7, M9 $\rightarrow$ \texttt{IMP-RULES-001}
$\rightarrow$ \texttt{TEST-RULES-M9-001}
$\rightarrow$ futuro arquivo de Rules $\rightarrow$ resultado ainda não produzido
\end{center}

Não se preenchem resultados futuros e não se cria rastreabilidade retroativa por
inferência. A evidência deve nascer da execução registrada.
